\chapter{Correlations, projections and degeneracies}\label{ch:fits}

\section{Projections}

The $\chi^2$ of Chapter~\ref{ch:chi2} tests one point of parameter space. To
measure one parameter, say $\theta_{13}$, all the others must be allowed to
take the values that fit the data best, since they are not known exactly and
can mimic a different $\theta_{13}$. This is the \emph{projection} of the
$\chi^2$ onto the $\theta_{13}$ axis (or \emph{profiling}, or
\emph{marginalisation}):
\begin{equation}
\chi^2(\theta_{13}) = \min_{\theta_{12},\theta_{23},\dcp,\dm{21},\dm{31},\rho}
  \Bigl[\chi^2(\theta_{12},\theta_{13},\dots,\rho) + \chi^2_\text{prior}\Bigr].
\label{eq:projection}
\end{equation}
More generally, any set of parameters can be kept fixed and the others
minimised over. The $\chi^2$ of a projection onto $k$ parameters is
distributed as $\chi^2$ with $k$ degrees of freedom; Table~\ref{tab:cl} gives
the values for the usual confidence levels.

\begin{table}[t]
\centering\small
\begin{tabular}{lccccc}
\toprule
& $1\sigma$ (68.27\%) & 90\% & $2\sigma$ (95.45\%) & 99\% & $3\sigma$ (99.73\%) \\
\midrule
1 d.o.f. & 1.00 & 2.71 & 4.00 & 6.63 & 9.00 \\
2 d.o.f. & 2.30 & 4.61 & 6.18 & 9.21 & 11.83 \\
3 d.o.f. & 3.53 & 6.25 & 8.02 & 11.34 & 14.16 \\
\bottomrule
\end{tabular}
\caption{$\Delta\chi^2$ for confidence levels and 1--3 parameters (degrees of
freedom).}
\label{tab:cl}
\end{table}

\section{External information: priors}\label{sec:priors}

Parameters that an experiment does not measure well (like the solar
parameters for a long-baseline experiment) or the matter density are
constrained by external information, as Gaussian priors:
\begin{equation}
\chi^2_\text{prior} = \sum_{k\ \text{free}}\frac{(x_k-\hat x_k)^2}{\sigma_k^2}
  + \sum_{e\ \text{free}}\frac{(\rho_e-\hat\rho_e)^2}{\sigma_\rho^2},
\label{eq:prior}
\end{equation}
summed over the fitted parameters. The centres $\hat x_k$ are the simulated
(true) values unless given otherwise; parameters without $\sigma_k$ have no
prior. Any other prior is added as a function of the parameters.

\section{Fitting}

\begin{apifun}{marginalize(start::OscParams; free, fixed, priors=(;), centers=simulated\_params(),\\
\phantom{marginalize(}prior=nothing, experiment=:all, setup=current\_setup()) -> FitResult}
Minimise $\chi^2+\chi^2_\text{prior}$ over the parameters listed in
\code{free} (or over all except those in \code{fixed}; by default over all),
starting at \code{start}, with the fixed parameters at their values in
\code{start}. Parameter names are those of Table~\ref{tab:params} and
\code{:ρ} for the matter densities of all experiments.
\code{priors} are the $\sigma_k$ as a named tuple, \code{centers} the
$\hat x_k$, and \code{prior} a function \code{OscParams -> Float64} added to
$\chi^2$. The result has the fields \code{chi2}, \code{params} (the parameters
at the minimum), \code{free} and \code{iterations} (the number of $\chi^2$
evaluations).
\end{apifun}

\begin{julia}
test = with_param(truth, :sin²2θ13, 0.08)
r = marginalize(test; fixed=(:θ13,),
                priors=(θ12=0.03truth.θ12, Δm21=0.03truth.Δm21, ρ=0.05))
r.chi2, r.params.δ, r.params.ρ

# only δ and the density free, a non-Gaussian prior on θ23
marginalize(test; free=(:δ, :ρ), priors=(ρ=0.05,),
            prior = p -> p.sin²θ23 < 0.5 ? 0.0 : 4.0)
\end{julia}

\begin{note}
The minimiser finds the minimum nearest to \code{start} (a local minimum).
Start fits at the true parameters, or, to find degeneracies, at the points
where they are expected (Section~\ref{sec:degeneracies}).
\end{note}

\section{Scans}

\begin{apifun}{chi2\_scan(base, x => xs; kwargs...) -> Vector\\
chi2\_scan(base, x => xs, y => ys; kwargs...) -> Matrix}
$\chi^2$ on a grid of one or two parameters, with all others as in
\code{base}. The parameter names include the derived ones of
Table~\ref{tab:params}, such as \code{:sin²2θ13}. Without keyword arguments,
each point is \code{chi2}; with any of the keywords of \code{marginalize},
each point is minimised over the other parameters, the scanned ones fixed.
The matrix has one row per value of \code{x}.
\end{apifun}

\noindent Figure~\ref{fig:delta} shows the projection onto $\dcp$ for NOvA,
Figure~\ref{fig:contours} a two-parameter scan for T2K and NOvA together:
\begin{julia}
δs = range(0, 2π, 73)
sys = chi2_scan(truth, :δ => δs)                                   # systematics only
fit = chi2_scan(truth, :δ => δs; free=(:θ23, :Δm31, :ρ), priors=(ρ=0.05,))
grid = chi2_scan(truth, :sin²2θ13 => range(0.06, 0.115, 56), :δ => δs)
\end{julia}

\begin{figure}[t]
\centering
\begin{tikzpicture}
\begin{axis}[wide,height=0.4\linewidth,xlabel={$\dcp/\pi$},ylabel={$\Delta\chi^2$},xmin=0,xmax=2,ymin=0,
  legend pos=north west,legend columns=1]
\addplot table[x=delta,y=sys]{figures/data/nova_delta.dat};
\addplot table[x=delta,y=fit]{figures/data/nova_delta.dat};
\addplot table[x=delta,y=fit_io]{figures/data/nova_delta.dat};
\addplot[gray,thin,dashed,domain=0:2] {1}; \addplot[gray,thin,dashed,domain=0:2] {4};
\addplot[gray,thin,dashed,domain=0:2] {9};
\legend{systematics only,fit of $\theta_{23}$\,$\dm{31}$\,$\rho$ (normal),fit started at inverted ordering}
\end{axis}
\end{tikzpicture}
\caption{$\Delta\chi^2$ as a function of the tested $\dcp$ for \file{NOvA.glb},
true $\dcp=3\pi/2$ and normal ordering: with all other parameters fixed at the
true values, minimised over $\theta_{23}$, $\dm{31}$ and $\rho$ (5\% prior),
and the same fit started at $\dm{31}\to-\dm{31}+\dm{21}$, which finds the
best fit with inverted ordering (the sign degeneracy). Dashed lines: $1$, $2$
and $3\sigma$ for one parameter.}
\label{fig:delta}
\end{figure}

\begin{figure}[t]
\centering
\begin{tikzpicture}
\begin{axis}[wide,height=0.45\linewidth,xlabel={$\sst{13}$},ylabel={$\dcp/\pi$},xmin=0.06,xmax=0.115,
  ymin=0,ymax=2,legend pos=north east,
  xticklabel style={/pgf/number format/fixed,/pgf/number format/precision=3},scaled x ticks=false]
\addplot[blue!70!black,thick] table[x=x,y=y]{figures/data/t2k_nova_1s.dat};
\addplot[red!70!black,thick] table[x=x,y=y]{figures/data/t2k_nova_2s.dat};
\addplot[green!50!black,thick] table[x=x,y=y]{figures/data/t2k_nova_3s.dat};
\addplot[only marks,mark=*,mark size=1.5pt,black] coordinates {(0.086064,1.5)};
\legend{$1\sigma$,$2\sigma$,$3\sigma$}
\end{axis}
\end{tikzpicture}
\caption{Allowed regions (2 degrees of freedom) in the $\sst{13}$--$\dcp$
plane for T2K and NOvA together, systematics only (\code{chi2\_scan} with two
parameters on a $56\times73$ grid). The dot marks the true values.}
\label{fig:contours}
\end{figure}

\section{The minimiser}\label{sec:minimizer}

GLoBES minimises with Powell's conjugate direction method and Brent's line
minimisation~\cite{NR}, which needs no derivatives. Two strategies are
available, chosen with \code{minimizer!}:
\begin{description}
\item[\code{:nested}] (default) Powell's method over the free oscillation
  parameters. At every step, the systematics of each rule are minimised
  separately (Eq.~\eqref{eq:pull}), each with its own Powell minimisation.
\item[\code{:hybrid}] One Powell minimisation over the free oscillation
  parameters and all nuisance parameters together, which can be faster with
  many rules.
\end{description}
The tolerances are $10^{-5}$ (relative) for oscillation parameters and
$10^{-8}$ for systematics. Since Powell's method can stop early in narrow
curved valleys, the minimisation over oscillation parameters is restarted with
fresh directions until it no longer improves (at most 10 times).

\begin{cdiff}
C GLoBES writes Powell's test for replacing a direction with a macro
\code{SQR(a)} that stores its argument in a static variable. Used twice in one
expression this is undefined behaviour, and GCC computes a different formula
than intended. C GLoBES also does not restart. Both can make fits stop at
points that are not minima (Section~\ref{sec:cdiff-powell}).
\end{cdiff}

\section{Degeneracies}\label{sec:degeneracies}

The appearance probability depends on $\theta_{13}$, $\dcp$, $\theta_{23}$ and
the sign of $\dm{31}$ in similar ways, so that different parameter sets can
give nearly the same event rates~\cite{Barger:2001yr}:
\begin{itemize}
\item the \emph{intrinsic} ($\theta_{13}$, $\dcp$) degeneracy;
\item the \emph{sign} degeneracy between normal and inverted ordering,
  ($\dm{31}\to-\dm{31}+\dm{21}$);
\item the \emph{octant} degeneracy, $\theta_{23}\to\pi/2-\theta_{23}$.
\end{itemize}
Each gives a local minimum of $\chi^2$. Since the minimiser only finds the
nearest minimum, a degenerate solution is found by starting the fit near it,
as in Figure~\ref{fig:delta}:
\begin{julia}
mirror = OscParams(truth; Δm31=-truth.Δm31 + truth.Δm21, θ23=π/2 - truth.θ23)
deg = marginalize(mirror; priors=(θ12=0.03truth.θ12, Δm21=0.03truth.Δm21, ρ=0.05))
deg.chi2, deg.params
\end{julia}
The projection including degeneracies is the smallest $\chi^2$ of all local
minima. A coarse scan of $\chi^2$ with systematics only (\code{chi2\_scan}) is
a good way to find where to start the fits.
